


\documentclass[12pt]{article}

\usepackage[margin = .8in]{geometry}
\usepackage{amsmath}
\usepackage{graphicx}
\usepackage{multicol, enumerate, tabularx}

\usepackage[parfill]{parskip}

\usepackage{adjustbox, soul}

\usepackage{fancyhdr}
\pagestyle{fancy}

\lhead{\emph{Math F113X Finance 1: Spreadsheets, Simple and Compound Interest}}
\rhead{\emph{lecture notes}}

\renewcommand{\emph}[1]{\textsf{\textbf{#1}}}
\usepackage{tikz}
\usetikzlibrary{calc,trees,positioning,arrows,fit,shapes,through, backgrounds}
\usetikzlibrary{patterns}

\usetikzlibrary{decorations.markings}
\usetikzlibrary{arrows}

\usepackage{pgfplots}

\usepackage{longtable}
\usepackage{tabularx}

\newcommand{\ds}{\displaystyle}
\newcommand{\ans}[1][1in]{\rule{#1}{.5pt}}

\newcommand{\points}[1]{(#1 points.)}		% Trying to be lazy.

\usepackage{array}
\newcolumntype{L}[1]{>{\raggedright\let\newline\\\arraybackslash\hspace{0pt}}m{#1}}
\newcolumntype{C}[1]{>{\centering\let\newline\\\arraybackslash\hspace{0pt}}m{#1}}
\newcolumntype{R}[1]{>{\raggedleft\let\newline\\\arraybackslash\hspace{0pt}}m{#1}}
\newcommand{\red}[1]{\textcolor{red}{#1}}

\newcommand{\be}{\begin{enumerate}}
\newcommand{\ee}{\end{enumerate}}


\begin{document}
Getting started with spreadsheets. 
\begin{enumerate}
\item  Basic Computation
	\be
	\item To add 3 + 4, enter \fbox{{\tt=3+4}}%\verb`=3+4`
	\item To subtract 100-76, enter \fbox{{\tt=100-76}}%\verb`=100-76`
	\item To multiply 4 times 18, enter \fbox{{\tt=4*18}}%\verb`=4*18`
	\item To divide 0.05 by 12, enter \fbox{{\tt=0.05/12}}%\verb`=0.05/12`
	\item To calculate $5^{25}$, enter \fbox{{\tt=5\char`\^25}}%\verb`=5^25`
	\ee

\item Compute an 18\% tip on a \$35.75 bill. 
	\be
	\item What is 18\% as a decimal? \ans
	\item What should we enter into the spreadsheet? \ans
	\ee

\item Make a tip calculator using \emph{cell references}.
	\be
	\item Open a new spreadsheet and label it {\tt TIP}.
	\item In cells A1, B1, C1, D1, type column headers: \\
	\fbox{{\tt\vphantom{Bp}Bill Amount}}%
	\fbox{{\tt\vphantom{Bp}Tip Percent}}%
	\fbox{{\tt\vphantom{Bp}Tip Paid}}%
	\fbox{{\tt\vphantom{Bp}Total}}
	\item In cells A2, B2, C2, D2, type: \\
	\fbox{{\tt\vphantom{Bp}$=35.75$}}%
	\fbox{{\tt\vphantom{Bp}$=0.18$}}%
	\fbox{{\tt\vphantom{Bp}$=A2 * B2$}}%
	\fbox{{\tt\vphantom{Bp}$=A2 + C2$}}
	\item What calculation did we do?\\
	\vspace{0.5in}
	\item Use the calculator to determine the Total if the Bill Amount is \$76.80 and the Tip Percent is 20\%.\\
	\vspace{0.5in}
	\ee
\item Keeping the Tip Percent at 20\%, we  now demonstrate the \emph{Fill Down} feature. 
	\be
	\item Enter {10} into cell \verb`A2` and enter  =\verb'A2'+10 into cell \verb`A3`
	\item Select \verb'A3' and drag down until you get to 110.
	\item Drag down cells \verb`B2` through \verb`D2`.
	\item How much is the tip on a \$110 meal? \ans 

	How much is the final bill? \ans

	How large does the Bill Amount need to be for the tip to be more than \$15? \ans
	\ee

\item Make a Simple Interest Calculator\\

Suppose \$1000 is invested in a savings account that earns 3.2\% interest on the \$1000 each year.
	\begin{enumerate}
	\item Open a new spreadsheet called {\tt Interest}.
	\item In cells A1, B1, C1, D1, E1 type column headers: \\
		\fbox{{\tt\vphantom{Bp}Interest Rate}}%
		\fbox{{\tt\vphantom{Bp}Principal}}%
		\fbox{{\tt\vphantom{Bp}Year}}%
		\fbox{{\tt\vphantom{Bp}Interest}}%
		\fbox{{\tt\vphantom{Bp}Total}} 
	\item Determine what should be in cells A2, B2, C2, D2, E2:  \\
		\fbox{{\tt\vphantom{\large Bp}\hspace{1in}}}%
		\fbox{{\tt\vphantom{\large Bp}\hspace{1in}}}%
		\fbox{{\tt\vphantom{\large Bp}\hspace{1in}}}%
		\fbox{{\tt\vphantom{\large Bp}\hspace{1in}}}%
		\fbox{{\tt\vphantom{\large Bp}\hspace{1in}}} 
	\item Determine what should be in cells A3, B3, C3, D3, E3:  \\
		\fbox{{\tt\vphantom{\large Bp}\hspace{1in}}}%
		\fbox{{\tt\vphantom{\large Bp}\hspace{1in}}}%
		\fbox{{\tt\vphantom{\large Bp}\hspace{1in}}}%
		\fbox{{\tt\vphantom{\large Bp}\hspace{1in}}}%
		\fbox{{\tt\vphantom{\large Bp}\hspace{1in}}} 
	\item Determine what should be in cells A4, B4, C4, D4, E4:  \\
		\fbox{{\tt\vphantom{\large Bp}\hspace{1in}}}%
		\fbox{{\tt\vphantom{\large Bp}\hspace{1in}}}%
		\fbox{{\tt\vphantom{\large Bp}\hspace{1in}}}%
		\fbox{{\tt\vphantom{\large Bp}\hspace{1in}}}%
		\fbox{{\tt\vphantom{\large Bp}\hspace{1in}}} 
	\item Use the pull down feature to determine the Total in the savings account in 10 years. Write down your steps. \\
	\vfill
	\item How long will it take for the savings account to double in value? Write down your steps. \\
	\vfill
	\item How can you calculate the Total in the savings account in 50 years \emph{without} using the pull down feature? Then check your calculation using the pull down feature. \\
	\vfill
	\item If $P$ is the principal, $r$ is the interest rate, $t$ is the number of years, and $A$ is the accumulated amount of money in the savings account, write a formula for $A$ and check your answer.
	\vfill
	\end{enumerate}
\newpage
\item How To Export the \emph{entire workbook} as a \emph{PDF}.
	\begin{enumerate}
	\item Save your work!
	\item In upper left menu, click \fbox{{\tt File}} to produce a dropdown menu.
	\item In dropdown menu, select \fbox{{\tt Download}}
	\item In right-side menu, select \fbox{{\tt PDF}}
	\end{enumerate}
	\vfill
\item Quick Summary of Finance Terms

	\be
	\item \emph{principal}


	\vfill	
	\item \emph{interest rate}


		\vfill
	\item \emph{interest}


		\vfill

	\item \emph{accumulated amount} or  \emph{end amount} or \emph{principal plus interest} or \emph{future value}
	
		\vfill	
	\item \emph{Formula for Simple Interest over Time}
	
		\vfill
	
	\item Comment on time $t$.\\
		\vfill
	\ee


\end{enumerate}
\end{document}


